% A two-column review in the layout of Nature's npj journals — made up, after
% "Fusion of medical imaging and electronic health records using deep
% learning" (npj Digital Medicine, 2020) — with what Reflow once read badly
% in one: no "Abstract" over the abstract, so the byline was never read and
% its numbers were taken for citations; a badge between each name and its
% numbers; the institutions, an equal-contribution note and the address to
% write to in a footnote on the first page; subsections headed in the
% text's own type; a table of studies set sideways on a page of its own,
% each study's cells running on to lines of their own; and an Author
% Contributions section that names the authors by their initials.
% Rebuilt with: pdflatex nature-review.tex (texlive-latex-recommended, -extra).
\documentclass[twocolumn,10pt]{article}
\usepackage[T1]{fontenc}
\usepackage[margin=0.75in,columnsep=0.3in]{geometry}
\usepackage{amsmath,booktabs,pifont,lscape,array,xcolor}
\usepackage[hidelinks]{hyperref}
\pagestyle{empty}
\setlength{\parindent}{1em}
\newcommand{\badge}{\hspace{0.3em}\textcolor{green!50!black}{\rule{0.6em}{0.6em}}\hspace{0.15em}}
\newcommand{\sub}[1]{\par\vspace{1.2\baselineskip}\noindent #1\par\noindent\ignorespaces}
\makeatletter
\renewcommand{\@maketitle}{%
  {\raggedright\sffamily\bfseries\small REVIEW ARTICLE\par}\vspace{0.4em}
  {\raggedright\LARGE Fusion of imaging and health records in Northfield\par}
  \vspace{0.8em}
  {\raggedright\sffamily Ada Lindqvist\badge$^{1,2,4}$\ding{41}, Tomas Okafor\badge$^{2,4}$ and Mira Castell$^{1,3}$\par}
  \vspace{1.2em}}
\makeatother
\begin{document}
\twocolumn[{\maketitle
\noindent\colorbox{gray!10}{\parbox{\dimexpr\textwidth-2\fboxsep}{\small Advancements in deep learning carry the potential to make significant contributions to healthcare, particularly in fields that use medical imaging for diagnosis and treatment decisions. Models for radiology consider only pixel values, without the clinical context that physicians rely on. We review the fusion of medical imaging with electronic health records, and give guidelines for researchers who want to apply it.}}\par
\vspace{1.5em}}]
\begingroup
\renewcommand{\thefootnote}{}
\footnotetext{\footnotesize\noindent$^{1}$Department of Biomedical Data Science, Northfield University, Northfield, USA. $^{2}$Center for Imaging, Northfield University, Northfield, USA. $^{3}$Department of Radiology, Eastmoor University, Eastmoor, USA. $^{4}$These authors contributed equally: Ada Lindqvist, Tomas Okafor. \ding{41}email: adal@northfield.example.org}
\endgroup
\noindent{\sffamily\bfseries\small INTRODUCTION}\par
\noindent The practice of modern medicine relies heavily on synthesis of information and data from multiple sources; this includes imaging pixel data, structured laboratory data, and unstructured narrative data. In a survey of radiologists, the majority stated that clinical information had a significant impact on interpretation\textsuperscript{1,2}. Deep learning in healthcare is proliferating due to the potential for automated systems to offload work from busy physicians\textsuperscript{3}.

The recent medical imaging literature shows a similar trend, where both health records and pixel data are leveraged in a fusion paradigm for solving complex tasks which cannot readily be tackled by a single modality, and this paragraph runs on to fill its column.
\sub{Early fusion}
The majority of the studies that remained after full-text screening used early fusion to join the multimodal input, concatenating imaging features with clinical ones before feeding them into a single model, as described in Table 1 below.
\sub{Late fusion}
Late fusion was used in three of the seventeen included studies. Each applied a different aggregation strategy to combine the predictions of separate models into one.

\noindent{\sffamily\bfseries\small AUTHOR CONTRIBUTIONS}\par
\noindent A.L. and T.O. are co-first authors who contributed equally to this study. Concept and design: A.L., T.O. and M.C. Study selection: A.L. and T.O. Supervision: M.C.

\clearpage
\begin{landscape}
\noindent{\sffamily\bfseries\small Table 1.}\hspace{0.5em}{\small Overview of studies included in the review.}\par\smallskip
\small
\newcolumntype{R}[1]{>{\raggedright\arraybackslash}p{#1}}
\begin{tabular}{@{}R{0.6in}R{0.4in}R{1in}R{1.3in}R{1.6in}R{0.7in}R{1.4in}@{}}
\toprule
Fusion strategy & Year & Author & Outcome & Input: non-imaging data & Number of samples & Model performance\\
\midrule
Early & 2017 & Thung et al. & Diagnosis of Alzheimer's disease & Patient data (age, sex, education) Genetic data (APOE4) & 805 & Fusion: 63.6\% Accuracy MRI: 58.0\% Accuracy\\\addlinespace[4pt]
Early & 2018 & An et al. & Glaucoma classification & Patient data (age, sex, spherical equivalent) & 163 & Fusion: 87.8\% Accuracy\\\addlinespace[4pt]
Late & 2018 & Reda et al. & Prostate cancer diagnosis & PSA blood test & 18 & Fusion: 94.4\% Accuracy MRI: 88.89\% Accuracy\\
\bottomrule
\end{tabular}
\end{landscape}
\end{document}
